\documentclass[10p]{article}
\usepackage[a4paper,left=1in,right=1in,top=1in,bottom=1in]{geometry}
\usepackage{comment}
\usepackage[colorlinks]{hyperref}
\usepackage{times}
\usepackage{amsmath}
\usepackage{eurosym}
\hypersetup{draft,linkcolor=blue,citecolor=blue,filecolor=black,urlcolor=blue}



\begin{document}

% \begin{center}
% \textbf{UNIVERSITA' DEGLI STUDI DI TORINO}
% \end{center}
%
%\begin{center}
%\textbf{DICHIARAZIONE SOSTITUTIVA DI CERTIFICAZIONE\\
% (art 46 del D.P.R. n.445 del 28/12/2000)}
%\end{center}
%
%\begin{center}
%\textbf{DICHIARAZIONE SOSTITUTIVA DELL'ATTO DI NOTORIETA'\\
%(artt. 19 e 47 del D.P.R. n.445 del 28/12/2000)}
%\end{center}
%
%\noindent
%Il sottoscritto CANTORE CRISTIANO \\
% Codice fiscale CNTCST81E05G942N nato a POTENZA prov. PZ il 05/05/1981\\
%  residente a  Londra, Regno Unito, indirizzo 4 Fosbrooke House, CAP. SW82XH, telefono +447983382586 \\
%consapevole delle sanzioni penali, nel caso di dichiarazioni non veritiere e falsità negli atti, richiamate dall’art. 76 del D.P.R. 445/2000 e dalle leggi speciali in materia\\
%
%\begin{center}
% \textbf{Dichiara}
%\end{center}
%
% \noindent
%  che il proprio curriculum relativo all’attivit\'{a} scientifica e didattica \`e il seguente:\\


%Decreto Rettore Universit\'{a} di Roma "La Sapienza" n. 2692/2021 del 19.10.2021


\begin{center}
{\LARGE Prof. Cristiano Cantore\\ Curriculum Vitae}\\
\end{center}
Sapienza Universit\`{a} di Roma \hfill Telefono Ufficio: \texttt{+390649766257}\\
Dipartimento di Economia e Diritto \hfill email: \href{mailto:cristiano.cantore@uniroma1.it}{\nolinkurl{cristiano.cantore@uniroma1.it}} \\
Via del Castro Laurenziano, 9, 6sto Piano Ala B, 00161 Roma  \hfill homepage: 
\href{http://www.cristianocantore.com}{www.cristianocantore.com}\\
Data e luogo di nascita:  May 5, 1981 Potenza, Italy \hfill Cellulare: \texttt{ +39(0)3924396187}\\\
Nazionalit\`{a}:  Italiana e Britannica \hfill pec: \href{mailto:cristianocantore@pec.it}{cristianocantore@pec.it}\\\\
\\
{\large \bf Posizioni Correnti} \\*[-.8pc]
\underline{\hspace{6in}} \\
-\emph{Marzo 2023-}: {\bf Professore Ordinario di Politica Economica presso La Sapienza Universi\`{a} di Roma} \\
-\emph{Settembre 2012-}: {\bf Fellow presso il Centre for International Macroeconomics, University of Surrey (Regno Unito)}\\
-\emph{Settembre 2018-}: {\bf Fellow presso il Centre for Macroeconomics, London School of Economics (Regno Unito)}\\
-\emph{Luglio 2024-}: {\bf Senior Fellow presso il Rimini Centre for Economic Analysis (Italia)}\\
-\emph{Luglio 2025-}: {\bf Associate Editor per Macroeconomic Dynamics}\\
\\
{\large \bf Istruzione e formazione universitaria} \\*[-.8pc]
\underline{\hspace{6in}} \\
-\emph{Settembre 2006-Luglio 2011}: {\bf Dottorato di Ricerca (Ph.D). in Economia, University of Kent (Regno Unito)} \\
Dichiarato equipollente ai sensi dell'art. 74 del DPR 382/80 con decreto ministeriale n.1805 del 7/5/2012 \\
Commissione -  Esaminatore Esterno: Prof. Martin Ellison (Oxford); Esaminatore Interno: Dr. Katsuyuki Shibayama (Kent); Nessuna correzione richiesta. \\
\\
-\emph{Anno accademico 2004-2005}: {\bf Master in Economia, Universitat Pompeu Fabra, Barcelona (Spagna)} \\
\\
-\emph{1999-2004}: {\bf Laurea (vecchio ordinamento) in Economia, Statistica e Scienze Sociali, conseguita il 7/5/2004 presso l'Universit\`{a} Commerciale L. Bocconi, Milano (voto 107/110).} Tesi: Gli effetti economici della transizione democratica in America Latina (Relatore Prof. Guido Tabellini).\\
\\
\\
{\large \bf Esperienze lavorative precedenti} \\*[-.8pc]
\underline{\hspace{6in}} \\
-\emph{Luglio 2021-Febbraio 2023}: {\bf Research Advisor presso il 'Research Hub' della Bank of England (Regno Unito)} \\
\\
-\emph{Febbraio 2021-Febbraio 2023}: {\bf Direttore del polo della Bank of England presso il Centre for Macroeconomics della London School of Economics (Regno Unito)} \\
\\
-\emph{Agosto 2018-Febbraio 2023}: {\bf 'Reader' in Scienze Economiche, presso University of Surrey, Guildford (Regno Unito)} \\
\\
-\emph{Gennaio 2022-Luglio 2022}: {\bf Professore a contratto, presso London School of Economics, Londra (Regno Unito)} \\
\\
-\emph{Settembre 2018-Giugno 2021}: {\bf Senior Research Economist presso il 'Research Hub' della Bank of England (Regno Unito)} \\
\\
-\emph{Aprile 2013-Luglio 2018}: {\bf Professore Associato (Senior Lecturer) in Scienze Economiche, presso University of Surrey, Guildford (Regno Unito)} \\
\\
-\emph{Ottobre 2021}: {\bf Visiting Professor presso il Collegio Carlo Alberto, Torino (Italia)} \\
\\
-\emph{Febbraio-Luglio 2014}: {\bf Visiting Professor presso la University of California San Diego (USA)} \\
\\
-\emph{Settembre 2009-Marzo 2013}: {\bf Ricercatore (Lecturer) in  Scienze Economiche, presso University of Surrey, Guildford (Regno Unito)} \\
\\
-\emph{Aprile-Agosto 2012}: {\bf Visiting Research Fellow presso il Banco de Espa\~{n}a, Madrid (Spagna)} \\
\\
-\emph{Gennaio 2012}: {\bf Visiting Professor presso la Facolt\`a di Economia, Universit\`a  di Cagliari} \\
\\
-\emph{Agosto-Ottobre 2008}: {\bf Tirocinio presso Banca Centrale Europea, DG Economia, Divisione "Sviluppo Macroeconomico Area Euro", Francoforte (Germania)}\\
\\
-\emph{Luglio-Settembre 2007}: {\bf Tirocinio presso OCSE, Dipartimento Economico, Assistente di ricerca per il report sulla situazione economica in Spagna del 2008, Parigi (Francia)}\\
\\
\\
{\large \bf Attivit\`a Didattica } \\*[-.8pc]
\underline{\hspace{6in}} \\
{\bf Einaudi Institute for Economics and Finance, Roma}\\
-\emph{Anno accademico \textbf{2026/2027}}: \\
Dottorato in Economia, titolare dei seguenti insegnamenti:\\
HANK models;\\
\\
{\bf Sapienza Universit\`{a} di Roma}\\
-\emph{Anno accademico \textbf{2026/2027}}: \\
Laurea Triennale, titolare dei seguenti insegnamenti:\\
Economia Politica;\\
Laurea Magistrale, titolare dei seguenti insegnamenti:\\
Macroeconomic Theory and Policy;\\
Laboratory on Macro Computation;\\
Relatore di 2 tesi di Laurea Triennale.\\
\\
{\bf ZUEL-SUR, Wuhan (Cina)}\\
-\emph{Anno accademico \textbf{2026/2027}}: \\
Laurea Magistrale, titolare dei seguenti insegnamenti:\\
Laboratory on Macro Computation;\\
\\
{\bf Einaudi Institute for Economics and Finance, Roma}\\
-\emph{Anno accademico \textbf{2025/2026}}: \\
Dottorato in Economia, titolare dei seguenti insegnamenti:\\
HANK models;\\
\\
{\bf Sapienza Universit\`{a} di Roma}\\
-\emph{Anno accademico \textbf{2025/2026}}: \\
Laurea Triennale, titolare dei seguenti insegnamenti:\\
Economia Politica;\\
Laurea Magistrale, titolare dei seguenti insegnamenti:\\
Macroeconomic Theory and Policy;\\
Relatore di 2 tesi di Laurea Triennale.\\
\\
{\bf Sapienza Universit\`{a} di Roma}\\
-\emph{Anno accademico \textbf{2024/2025}}: \\
Laurea Triennale, titolare dei seguenti insegnamenti:\\
Politica Economica;\\
Laurea Magistrale, titolare dei seguenti insegnamenti:\\
Teorie e politiche Macroeconomiche;\\
Relatore di 2 tesi di Laurea Magistrale.\\
\\
{\bf Sapienza Universit\`{a} di Roma}\\
-\emph{Anno accademico \textbf{2023/2024}}: \\
Laurea Triennale, titolare dei seguenti insegnamenti:\\
Politica Economica;\\
Laurea Magistrale, titolare dei seguenti insegnamenti:\\
Teorie e politiche Macroeconomiche;\\
Relatore di 1 tesi di Laurea Magistrale.\\
\\
{\bf Sapienza Universit\`{a} di Roma}\\
-\emph{Anno accademico \textbf{2022/2023}}: \\
Laurea Magistrale, titolare dei seguenti insegnamenti:\\
Economic Policy;\\
\\
{\bf University of Surrey, Guildford (Regno Unito)}\\
-\emph{Anno accademico \textbf{2022/2023}}: \\
Corso di Dottorato di Ricerca (PhD), titolare dei seguenti insegnamenti:\\
Macroeconomia Avanzata I (Advanced Macroeconomics I);\\
Relatore di 1 tesi di dottorato di ricerca.\\
Corsi di laurea di primo livello (underaduate), titolare dei seguenti insegnamenti: \\
Questioni contemporanee in Economia (Contemporary Isssues in Economics);\\
\\
{\bf London School of Economics, Londra (Regno Unito)}\\
-\emph{Anno accademico \textbf{2021/2022}}: \\
Corso di Master, titolare dei seguenti insegnamenti:\\
Economia Monetaria e fluttuazioni economiche (Monetary Economics and aggregate fluctuations);\\
\\
{\bf Collegio Carlo Alberto, Torino (Italia)}\\
-\emph{Anno accademico \textbf{2021/2022}}: \\
Corso di Master, titolare dei seguenti insegnamenti:\\
Introduzione alla Macroeconomia Dinamica (Intro to Dynamic Macroeconomics);\\
\\
{\bf University of Surrey, Guildford (Regno Unito)}\\
-\emph{Anno accademico \textbf{2021/2022}}: \\
Corso di Dottorato di Ricerca (PhD), titolare dei seguenti insegnamenti:\\
Macroeconomia Avanzata I (Advanced Macroeconomics I);\\
Relatore di 1 tesi di dottorato di ricerca.\\
\\
-\emph{Anno accademico \textbf{2020/2021}}: \\
Corso di Dottorato di Ricerca (PhD), titolare dei seguenti insegnamenti:\\
Macroeconomia Avanzata I (Advanced Macroeconomics I);\\
Corsi di laurea di primo livello (undergraduate), titolare dei seguenti insegnamenti: \\
Macroeconomia Avanzata (Topics in Macroeconomics).\\
Relatore di 1 tesi di dottorato di ricerca.\\
\\
-\emph{Anno accademico \textbf{2019/2020}}: \\
Corso di Dottorato di Ricerca (PhD), titolare dei seguenti insegnamenti:\\
Macroeconomia Avanzata I (Advanced Macroeconomics I);\\
Macroeconomia Applicata (Topics in Empirical Macroeconomics).\\
Relatore di 1 tesi di dottorato di ricerca.\\
\\
-\emph{Anno accademico \textbf{2018/2019}}: \\
Corso di Dottorato di Ricerca (PhD), titolare dei seguenti insegnamenti:\\
Macroeconomia Avanzata I (Advanced Macroeconomics I);\\
Macroeconomia Applicata (Topics in Empirical Macroeconomics).\\
Corsi di laurea di secondo livello (postgraduate), titolare dei seguenti insegnamenti: \\
Macroeconomia (Macroeconomics).\\
Relatore di 3 tesi di dottorato di ricerca.\\
\\
-\emph{Anno accademico \textbf{2017/2018}}: \\
Corso di Dottorato di Ricerca (PhD), titolare dei seguenti insegnamenti:\\
Macroeconomia Avanzata I (Advanced Macroeconomics I);\\
Macroeconomia Applicata (Topics in Empirical Macroeconomics).\\
Relatore di 3 tesi di laurea magristrale e di 3 tesi di dottorato di ricerca.\\
Tutor di 35 studenti iscritti ai corsi di laurea di primo livello e di master del dipartimento di economia.\\
\\
-\emph{Anni accademici \textbf{2015/2016} e \textbf{2016/2017}}: \\
Corsi di laurea di primo livello (undergraduate), titolare dei seguenti insegnamenti: \\
Teoria Macroeconomica (Intermediate Macroeconomics 2); \\
Corso di Dottorato di Ricerca (PhD), titolare dei seguenti insegnamenti:\\
Macroeconomia Avanzata I (Advanced Macroeconomics I);\\
Macroeconomia Applicata (Topics in Empirical Macroeconomics).\\
Relatore di 3 tesi di laurea magristrale e di 3 tesi di dottorato di ricerca.\\
Tutor di 25 studenti iscritti ai corsi di laurea di primo livello e di master del dipartimento di economia.\\
\\
-\emph{Anni accademici \textbf{2013/2014} e \textbf{2014/2015}}: \\
Corsi di laurea di primo livello (undergraduate), titolare dei seguenti insegnamenti: \\
Teoria Macroeconomica (Intermediate Macroeconomics 2); \\
Analisi dei dati per l'economia (Economic Data Analysis).\\
Corsi di laurea di secondo livello (postgraduate), titolare dei seguenti insegnamenti: \\
Commercio Internazionale (Microeconomics of the Global Economy).\\
Relatore di 4 tesi di laurea magristrale e di 3  tesi di dottorato di ricerca.\\
Tutor di 25 studenti iscritti ai corsi di laurea di primo livello e di master del dipartimento di economia.\\
\\
-\emph{Anni accademici \textbf{2011/2012} e \textbf{2012/2013}}: \\
Corsi di laurea di primo livello (undergraduate), titolare dei seguenti insegnamenti: \\
Teoria Macroeconomica (Intermediate Macroeconomics 2); \\
Analisi dei dati per l'economia (Economic Data Analysis).\\
Corsi di laurea di secondo livello (postgraduate), titolare dei seguenti insegnamenti: \\
Teoria Economica Avanzata (Advanced Economic Theory).\\
Commercio Internazionale (Microeconomics of the Global Economy). \\
Relatore di 3 tesi di laurea magristrale e di 3 tesi di dottorato di ricerca.\\
Tutor di 20 studenti iscritti ai corsi di laurea di primo livello e di master del dipartimento di economia.\\
\\
-\emph{Anno accademico \textbf{2010/2011}}: \\
Corsi di laurea di primo livello (undergraduate), titolare dei seguenti insegnamenti: \\
Teoria Macroeconomica (Intermediate Macroeconomics 2); \\
Economia e Finanza Internazionale (International Economics and Finance);\\
Corsi di laurea di secondo livello (postgraduate), titolare dei seguenti insegnamenti: \\
Teoria Economica Avanzata (Advanced Economic Theory).\\
Relatore di 3 tesi di laurea magristrale e di 1 tesi di dottorato di ricerca.\\
Tutor di 20 studenti iscritti ai corsi di laurea di primo livello e di master del dipartimento di economia.\\
\\
-\emph{Anno accademico \textbf{2009/2010}}: \\
Corsi di laurea di primo livello (undergraduate), titolare dei seguenti insegnamenti: \\
Teoria Macroeconomica (Intermediate Macroeconomics Theory); \\
Economia e Finanza Internazionale (International Economics and Finance).\\
Relatore di 3 tesi di laurea magristrale.\\
Tutor di 20 studenti iscritti ai corsi di laurea di primo livello e di master del dipartimento di economia.\\
\\
{\bf University of Kent, Canterbury (Regno Unito)}\\
-\emph{Anni accademici \textbf{2006/2007}, \textbf{2007/2008} e \textbf{2008/2009}}: \\
Corsi di laurea di primo livello (undergraduate), esercitazioni per i seguenti insegnamenti: \\
Matematica per l'economia (Mathematics for Economics);\\
Statistica per l'Economia (Statistics for Economics);\\
Economia Europea Contemporanea (Contemporaneous European Economics);
Finanza Internazionale (International Finance);
\\
\\
{\bf \underline{Supervisione e valutazione di studenti che hanno terminato il dottorato di ricerca in economia}}: \\
Relatore: Federico Pilla (Univeristy of Surrey - Regno Unito, 2026); Vedanta Dhamija (Univeristy of Surrey - Regno Unito, 2023); Joseph Grilli (Univeristy of Surrey - Regno Unito, 2019);  Tae Min Lee (Univeristy of Surrey - Regno Unito, 2017); Marcelo Almeida (Univeristy of Surrey - Regno Unito, 2017); Matteo Ghilardi (Univeristy of Surrey - Regno Unito, 2013).\\
Componente della commissione dell' esame finale: Frantisek Masek (Sapienza Universit\`{a} di Roma - Italia, 2026); Gregorio Ghetti (Università di Milano Bicocca - 2025); Giuliano de Queiroz Ferreira (Universidade de S\~ao Paulo - Brasile, 2023); Silvio Tunis (Universit\`{a} di Cagliari - Italia, 2023); Alessandro Cantelmo (London City University - Regno Unito, 2018); Fred Iklaga (Univeristy of Surrey - Regno Unito, 2017); Diana Lima (Univeristy of Surrey - Regno Unito, 2017); Jonathan Swarbrick (Univeristy of Surrey - Regno Unito, 2016); Veronica Arcurio Vasconez (Universit\`a Paris 1 - Francia, 2015).
\\
\\
{\bf \underline{Incarichi di docenza per corsi di formazione professionale e Summer Schools}} \\
\emph{17 Settembre 2025:} {\bf CIMS Summer School, University of Surrey (Regno Unito).}\\
\emph{24-26 Febbraio 2025:} {\bf Banca Centrale del Paraguay, Assunción.}\\
\emph{9 Dicembre 2024:} {\bf CIMS Winter School, University of Surrey (Regno Unito).}\\
\emph{10 Settembre 2024:} {\bf CIMS Summer School, University of Surrey (Regno Unito).}\\
\emph{15-26 Marzo \& 4-5 Aprile 2024:} {\bf Banca Centrale del Paraguay, Online.}\\
\emph{11-15 Settembre 2023:} {\bf CIMS Summer School, University of Surrey (Regno Unito).}\\
\emph{12-16 Settembre 2022:} {\bf CIMS Summer School, University of Surrey (Regno Unito).}\\
\emph{7-14 Settembre 2021:} {\bf CIMS Summer School, University of Surrey (Regno Unito).}\\
\emph{7-11 Settembre 2020:} {\bf CIMS Summer School, University of Surrey (Regno Unito).}\\
\emph{9-13 Settembre 2019:} {\bf CIMS Summer School, University of Surrey (Regno Unito).}\\
\emph{3-7 Settembre 2018:} {\bf CIMS Summer School, University of Surrey (Regno Unito).}\\
\emph{4-8 Settembre 2017:} {\bf CIMS Summer School, University of Surrey (Regno Unito).}\\
\emph{25-26 Luglio 2017:} {\bf DIW, Berlino (Germania).}\\
\emph{30 Agosto-3 Settembre 2016:} {\bf CIMS Summer School, University of Surrey (Regno Unito).}\\
\emph{6-10 Giugno 2016:} {\bf Central Bank of Nigeria, Abuja (Nigeria).}\\
\emph{25-29 Gennaio 2016:} {\bf Central Bank of Nigeria, Abuja (Nigeria).}\\
\emph{7-11 Settembre 2015:} {\bf CIMS Summer School, University of Surrey (Regno Unito).}\\
\emph{6-9 Luglio 2015:} {\bf Computational Economics Summer School, Universit\`{a} di Cagliari.}\\
\emph{20-22 Aprile 2015:} {\bf CIMS Easter School, University of Surrey (Regno Unito).}\\
\emph{8-13 Settembre2014:} {\bf CIMS Summer School, University of Surrey (Regno Unito).}\\
\emph{9-13 Settembre 2013:} {\bf CIMS Summer School, University of Surrey (Regno Unito).}\\
\emph{14-18 Aprile 2013:}{\bf RES Easter School, University of Birmingham (Regno Unito).}\\
\emph{1014 Settembre 2012:} {\bf CIMS Summer School, University of Surrey(Regno Unito).}\\
\emph{30 Gennaio-2 Febbraio 2012:} {\bf PhD program, University of Porto (Portogallo).}\\
\emph{5-6 Luglio 2011:} {\bf Scottish PhD program, University of Glasgow (Regno Unito).}\\
\emph{31 Gennaio-4 Febbraio 2011:} {\bf National Institute of Public Policy and Finance, New Delhi (India).}\\
\\
\\
{\large \bf Attivit\`{a} di Ricerca} \\*[-.8pc]
\underline{\hspace{6in}} \\
\\
{\bf \underline{Pubblicazioni}}\\
-(2026) \emph{A tail of labor supply and a tale of monetary policy}, joint with Filippo Ferroni (University of Bologna), Haroon Mumtaz (Queen Mary University of London) and Angeliki Theophilopoulou (Brunel University London).
\textbf{Journal of the European Economic Association}, jvag030.\\
\\
-(2025) \emph{Monetary–fiscal interaction and the liquidity of government debt}, joint with Edoardo Leonardi (London School of Economics).
\textbf{European Economic Review}, 173, 104979.\\
\\
-(2025) \emph{\href{https://www.elgaronline.com/edcollchap/book/9781035329397/chapter13.xml}{The economics of intellectual property protection for cannabis: challenges, opportunities and the evolution of patents}}. in '\textbf{Intellectual Property and Cannabis}', edito da N. Corthésy, E. Bonadio e N. Williams, Edward Elgar Publishing.\\
\\
-(2024) \emph{Unwinding quantitative easing: State dependency and household heterogeneity}, joint with Pascal Meichtry (Banque de France).
\textbf{European Economic Review}, 170, 104865.\\
\\
-(2021) \emph{Workers, Capitalists, and the Government: Fiscal Policy and Income (Re)Distribution}, joint with Lukas Freund (University of Cambridge), \textbf{Jounrnal of Monetary Economics}, 119, 58-74.\\
\\
-2020 \emph{The Missing Link: Labor Share and Monetary Policy}, joint with Filippo Ferroni (Chicago FED) and Miguel Le\'{o}n-Ledesma (University of Kent).
\textbf{Journal of the European Economic Association}, \\ https://doi.org/10.1093/jeea/jvaa034, \textbf{in stampa}.\\
\\
-2019: \emph{Optimal Fiscal and Monetary Policy, Debt Crisis and
Management}, con  Paul Levine (University of Surrey), Giovanni Melina (City University London) and Joseph Pearlman (City University London). \textbf{Macroeconomic Dynamics}, 23-3, 1166-1204. \\
\\
-2017: \emph{The Dynamics of hours worked and technology}, con Filippo Ferroni (Chicago FED) and Miguel Le\'{o}n-Ledesma (University of Kent).
\textbf{Journal of Economic Dynamics and Control}, 82, 67-82.\\
\\
-2015: \emph{CES technology and Business Cycle fluctuations}, con Paul Levine (University of Surrey), Joseph Pearlman (City University London) e Bo Yang (Xi'an Jiaotong - Liverpool University).
\textbf{Journal of Economic Dynamics and Control}, 61(C), 133-151.\\
\\
-2014: \emph{Shocking Stuff: Technology, Hours and Factor Substitution}, con Miguel Le\'{o}n-Ledesma (University of Kent), Peter McAdam (ECB) e Alpo Willman (ECB). 
\textbf{Journal of the European Economic Association}, 12-1, 108-128.\\
\\
-2014: \emph{A Fiscal Stimulus and Jobless Recovery}, con  Paul Levine (University of Surrey) e Giovanni Melina (IMF). \textbf{Scandinavian Journal of Economics}, 116-3, 669-701.\\
\\
-2013: \emph{The Science and  Art of DSGE Modelling I. Construction and Bayesian Estimation}, con V. Gabriel,  P. Levine, J. Pearlman e B. Yang.   in  '\textbf{Handbook of Research Methods and Applications}', edito da N. Hashimzade and M. Thornton,  Edward Elgar Publishing.\\
\\
-2013: \emph{The Science and  Art of DSGE Modelling  II.~Model Comparisons, Model Validation, Policy Analysis and General Discussion}, con V. Gabriel,  P. Levine, J. Pearlman e B. Yang. in  '\textbf{Handbook of Research Methods and Applications}', edito da N. Hashimzade and M. Thornton,  Edward Elgar Publishing.\\
\\
-2012: \emph{A Fiscal Stimulus with Deep Habits and Optimal Monetary Policy}, con Paul Levine (University of Surrey), Giovanni Melina (IMF) e Bo Yang (University of Surrey). \textbf{Economics Letters}, 117-1, 348-353.\\
\\
-2012: \emph{Getting Normalization Right: Dealing with 'Dimensional Constants' in Macroeconomics}, e Paul Levine (University of Surrey). \textbf{Journal of Economic Dynamics and Control},  36, 1931-1949.\\
\\
\\
{\bf \underline{Working papers}} \\
-2018: \emph{Deep versus Superficial Habits: It's all in the persistence}, con  Paul Levine (University of Surrey), Giovanni Melina (City University London). \textbf{working paper permanente}\\
\\
\\
{\bf \underline{Borse di studio, fondi di ricerca, premi e partecipazioni a comitati scientifici}}  \\
-\emph{2026-}: {\bf Organizzatore conferenza 3rd Sapienza Macro Meetings, Roma (Italia)}.\\
\\
-\emph{2026-}: {\bf Organizzatore conferenza 6th Sailing the Macro Workshop, Ortygia (Italia)}.\\
\\
-\emph{2025-}: {\bf Organizzatore conferenza 2nd Sapienza Macro Meetings, Roma (Italia)}.\\
\\
-\emph{2025-}: {\bf Organizzatore conferenza 5th Sailing the Macro Workshop, Ortygia (Italia)}.\\
\\
-\emph{2024-}: {\bf Organizzatore conferenza 1st Sapienza Macro Meetings, Roma (Italia)}.\\
\\
-\emph{2024-}: {\bf Organizzatore conferenza 4th Sailing the Macro Workshop, Ortygia (Italia)}.\\
\\
-\emph{2024-}: {\bf Organizzatore conferenza  Household Heterogeneity, Inflation, and the Green Transition, Sapienza University. \hfill 2024.
, Roma (Italia)}.\\
\\
-\emph{2024-}: {\bf Membro del ESRC Peer Review College. (Regno Unito)}.\\
\\
-\emph{2023-}: {\bf Organizzatore conferenza 3rd Sailing the Macro Workshop, Ortygia (Italia)}.\\
\\
-\emph{2023-2024}: {\bf Membro del comitato scientifico della conferenza annuale dellla Money, Macro and Finance Society (Regno Unito)}.\\
\\
-\emph{2023-2024}: {\bf Membro del comitato scientifico della Money, Macro and Finance Society (Regno Unito)}.\\
\\
-\emph{2018-2023}: {\bf Editore degli Staff Working Paper della Banca D'Inghilterra.}\\
\\
-\emph{Aprile 2022} {\bf Organizzatore conferenza \emph{\href{https://www.mmf.ac.uk/phd/2022-phd-conference/}{9th Annual MMF PhD Conference}}.}\\
\\
-\emph{2022} {\bf Organizzatore conferenza \emph{\href{https://www.bankofengland.co.uk/events/2022/february/first-annual-bear-conference}{First Annual BEAR Conference: The Monetary Toolkit}}.}\\
\\
-\emph{2021} {\bf Organizzatore conferenza \emph{\href{https://www.mmf.ac.uk/productivity-workshop/}{Productivity and Structural Change}} sponsorizzata da NIESR, The productivity institute, MMF e la Banca D'Inghilterra.}\\
\\
-\emph{2021} {\bf Organizzatore conferenza  \emph{\href{https://www.bankofengland.co.uk/events/2021/may/mmf-workshop-on-macroeconomic-consequences-of-technological-change}{Macroeconomic Consequences of Technological Change workshop}} sponsorizzata da MMF, University of Kent e la Banca D'Inghilterra.}\\
\\
-\emph{2020}: {\bf Organizzatore conferenza \emph{Income Distribution, Wealth Distribution and Central Bank Policies} sponsorizzata dalla London School of Economics e la Banca D'Inghilterra.}\\
\\
-\emph{2018}: {\bf Componente del comitato scientifico della Money, Macroeconomics and Finance PhD conference.}\\
\\
-\emph{2017-2018}: {\bf Fondo di ricerca della Commissione Europea: CO-CREATION-08-2016/2017 -  Modelling and evaluating the socio-economic impacts of research and innovation with the suite of macro- and regional-economic models.} 323,440\euro$\,$. Co-investigator \\
\\
-\emph{2013-2015}: {\bf Componente del comitato scientifico del Simposio de la Associaci\'on Española de Economia.}\\
\\
-\emph{Luglio 2012}: {\bf Royal Economic Society Special Project Grant per finanziare il \emph{2nd Workshop on Structural Change \& Macroeconomic Dynamics.}} 2,600\textsterling\\
\\
-\emph{Aprile 2012 - Agosto 2012}: {\bf Banco de Espa\~{n}a Fellowship 'Trabajos de investigaci\'{o}n sobre econom\'{i}a'.}\\
\\
-\emph{Ottobre 2010 - Settembre 2012}: {\bf Fondo di Ricerca ESRC grant: Monetary and Fiscal Policy Rules with Labour Market and Financial Frictions.} 347,000\textsterling$\,$. Co-investigator\\
\\
-\emph{Settembre 2006 - Settembre 2009}: {\bf Borsa di studio per partecipare al dottorato di ricerca in Economia, ESRC e  University of Kent (Regno Unito).}\\
\\
-\emph{Settembre 2009 - Settembre 2018}: {\bf Vice direttore del centro di ricerca 'Centre for International Macroeconomic Studies' (CIMS) presso l'University of Surrey (Regno Unito).}\\
\\
\\
{\bf \underline{Conferenze e seminari}:}\\
\textbf{2026:} {2nd Banca d'Italia Annual Research Conference
on Monetary Policy; University of the West Indies, Kingston; Bank of Jamaica, Kingston; EIEF, Roma; Università di Tor Vergata, Roma.}
\\
\textbf{2025:} {24th Macroeconomic Dynamics Workshop, Università LUISS Roma; PSE Macro Days 2025, Paris School of Economics, Francia; BSE Summer Forum, Barcelona, Spagna; Università di Bonn, Germania; Banca Centrale del Paraguay, Assunción.}
\\
\textbf{2024:} {2024 Quantitative Macroeconomics Workshop, Banca Centrale Australiana, Sidney; Banca Centrale della Colombia, online; 3rd Dynare Workshop for Advanced users, JRC Ispra (Italia); ICMAIF 2024, Universit\`{a} di Creta (Grecia); Padova Macro Talks, Universit\`{a} di Padova; Paris School of Economics, Paris; EIEF, Roma; Commissione Europea, DG ECFIN (Online); Society of Nonlinear Dynamics and Econometrics, Universit\`{a} di Padova.}
\\
\textbf{2023:} {ASSA Meetings, New Orleans (USA); University of York (Regno Unito); Bank of England (Regno Unito); Norges Bank (Norvegia); FINPRO 3, Banca d'Italia; EEA-ESEM23, Universitat Pompeu Fabra (Spagna), Workshop on Macroeconomics and Survey Data, Banca Nazionale del Belgio.}
\\
\textbf{2022:} {Universitat Autonoma de Barcelona (Spagna); Theories and Methods in Macroeconomics 2022, Kings College Londra (Regno Unito); Royal Economic Society (invited for policy discussion), University of Reading (Regno Unito); Banca d'Italia, online; Society for Nonlinear Dynamics and Econometrics, online; Barcelona Summer Forum: Monetary Policy and Central Banking, Universitat Autonoma de Barcelona (Spagna); Salento Macro Meetings, Lecce; Universit\`{a} La Sapienza, Roma; Universit\'{e} catholique de Louvain (Belgio); Household Heterogeneity conference, Banca di Francia, Parigi; 4th European Midwest Micro/Macro Conference (EM3C), Frankfurt School of Finance and Management, Francoforte.}
\\
\textbf{2021:} {11th ifo Conference on Macroeconomics and Survey Data, Ifo Institute; University of Surrey; CFM workshop, London School of Economics; 3TM, Banca di Francia; 8th Bank of Italy-CEPR Conference on Money, Banking and Finance: “Closing the Gaps: The Future of Stabilisation Policies after the COVID-19 Pandemic”, Banca d'Italia; EEA/ESEM 2021; 52nd MMF Annual conference, 1st Sailing the Macro Workshop, Ventotene; Collegio Carlo Alberto, Torino; Macroeconometric Workshop, King's College Londra; Universidad de Antioquia, Medell\'{i}n (Colombia); University of Surrey (Regno Unito); Workshop Empirical Monetary Economics 2021, SciencePo, Parigi (Francia); European Winter Meeting of the Econometric Society 2021, Universitat de Barcelona (Spagna).}
\\
\textbf{2020:} {Birkbeck, Londra (Regno Unito); EEA/ESEM 2020, University of Rotterdam (Olanda); 2020 Central Bank Macro Modelling Workshop, Norges Bank (Norvegia); XXV meeting del Central Bank Researchers Network (CEMLA), Banca Centrale dell' Uruguay.}
\\
\textbf{2019:} {Banca d'Italia, Roma; 7th Belgian Macroeonomic Wokshop (\textbf{invited lecture}), University of Ghent (Belgio); EEA/ESEM 2019, University of Manchester (Regno Unito); 3rd Workshop on 'Macroeconomic and Financial Time Series Analysis', University of Lancaster (Regno Unito); CEPR Macroeconomic Modelling and Model Comparison, Goethe University Francoforte (Germania); 2019 Royal Economic Society Annual Conference, University of Warwick (Regno Unito); International Monetary Fund, Washington DC (USA); 27th Annual Symposium of the Society for Nonlinear Dynamics and Econometrics, Dallas FED (USA); 'Labor Income Share Developments: Drivers and Policy Implications', Bank of Greece, Athens (Grecia).}
\\
\textbf{2018:} {Ridge Forum December 2018: Macroeconomics and Development, Buenos Aires (Argentina); 6th Workshop in Macro Banking and Finance, Alghero; 71st European Meeting of the Econometric Society, Colonia (Germania); Banca d'Inghilterra, Londra (Regno Unito); 14th Dynare Conference, Banca Centrale Europea, Francoforte (Germania); Deutsche Bundesbank, Francoforte (Germania); Universitat de Barcelona, Barcelona (Spagna); BETA - Universit\'{e} de Lorraine, Nancy (Francia).} 
\\
\textbf{2017:} {16th workshop "Macronomic Dynamics: theory and applications", Universit\`{a} Cattolica de Sacro Cuore Milano;  XI Workshop on Public Design: "Structural Changes, Labor Markets and Policy", Universitat de Girona (Spagna); Banca del Giappone, Tokyo (Giappone); 13th Dynare Conference, Universit\'{a} di Tokyo (Giappone); Universit\`{a} di Verona; DIW, Berlino (Germania); 1st Research Conference of the CEPR Network on Macroeconomic Modelling and Model Comparison MMCN, Goethe University, Francoforte (Germania); Banca d'Inghilterra, Londra (Regno Unito); Fondo Monetario Internazionale, Washington DC (USA); Federal Reserve Bank, Chicago (USA); Collegio Carlo Alberto, Torino; University of Surrey, Guildford (UK); University of Kent, Canterbury (UK).} 
\\
\textbf{2016:} {University of Portsmouth; Workshop '\emph{Developments in Macroeconomics}' Universit\`{a} di Roma La Sapienza; Lacea-Lames, Eafit, Medell\'{i}n; Spanish Economic Association Simposio, Bilbao.} 
\\
\textbf{2015:} {University of Glasgow, Business School; Mid-West Macro Meeting, University of Washington St. Louis; University of Manchester; Lacea meeting, Santa Cruz, Bolivia; Spanish Economic Association Simposio, Girona.} 
\\
\textbf{2014:} {Rice University, Huston (TX); University of California San Diego; $10^{th}$ Dynare conference, Banque de France, Paris; Lacea Lames meeting, Sao Paulo; Spanish Economic Association Simposio, Palma de Mallorca.} 
\\
\textbf{2013:} {Spanish Economic Association Simposio, Santander; Xi'an Jiaotong-Liverpool University, SuZhou China; 9th Dynare conference, SUFE Shanghai; MBF Workshop, Universita' di Milano Bicocca; Reserve Bank of New Zealand; Universita' di Pavia; Utrecht School of Economics; Universita' di Padova; Workshop on Empirical Macroeconomics, University of Ghent.}
\\
\textbf{2012:} {Spanish Economic Association Simposio, Vigo; Cardiff Business School; SIE, Matera; 8th Dynare Conference, Zurich; CIMS Conference, Guildford; EEA-ESEM University of Malaga; Structural Change and Macroeconomic Dynamics University of Cagliari; University of Porto; Banco de Espa\~{n}a; Workshop on Structural Change \& Macroeconomic Dynamics, University of Cagliari.}
\\
\textbf{2011:} {NIPFP, New Delhi; Structural Change and Macroeconomic Performance workshop at University of 
Cagliari.}
\\
\textbf{2010:} {MONFISPOL Conference, London Metropolitan University, London; Computational Economics and Finance 2010 Conference, Department of Economics,  City University, London; Spanish Economic Association Simposio, Department of Economics,  Universidad Autonoma de Madrid. }
\\
\textbf{2009:} {Spanish Economic Association Simposio, Department of Economics,  University of Valencia; Association of Southern European Economist Annual conference, Department of Economics, Bo\u{g}azi\c{c}i University, Istanbul; 41st Annual Conference, Money Macro and Finance research group, Bradford School of Management, Bradford University; Annual international conference on macroeconomic analysis and international finance, Department of Economics, University of Crete, Rethymno; Department of Financial Economics, University of Hannover; Department of Economics, University of Konstanz.}
\\
\textbf{2008:} {Macroeconomic and Financial Linkages: Theory and Practice. Peterhouse College, University of Cambridge; Association of Southern European Economic Theorists (ASSET) 2008 meeting. EUI, Florence; Doctoral Workshop on Dynamic Macroeconomics, Universite Louis Pasteur Strasbourg; Southern Workshop in Macroeconomics, University of Auckland and Reserve Bank of New Zealand.}
\\
\textbf{2007:} {Second Doctoral Workshop, EBIM, University of Bielefeld.}
\\
\\
{\large \bf Attivit\`{a} di consulenza} \\*[-.8pc]
\underline{\hspace{6in}} \\
-\emph{Giugno 2026-Febbraio 2027}: {\bf Consulente per la Asian Development Bank.}\\
\\
-\emph{Giugno 2024-Settembre 2025}: {\bf ``Impact of monetary policy on households' desired labour supply in the euro area''}, gara d'appalto DG-ECFIN, Commissione Europea.\\
\\
-\emph{Dicembre 2024-Gennaio 2026}: {\bf ``Methodology and measurement update of the Gender Equality Index''}, FWC EIGE/2023/OPER/05 - Lot 2. Consulenza per Fondazione Brodolini.\\
\\
\\
{\large \bf Incarichi Amministrativi presso University of Surrey, Guildford (Regno Unito)} \\*[-.8pc]
\underline{\hspace{6in}} \\
\emph{2026-}: Coordinatore del Dottorato di Ricerca in Economia Politica, Sapienza University of Rome.\\
\emph{2025-}: Presidente della commissione paritetica docenti e studenti della Facoltà di Economia, Sapienza University of Rome.\\
\emph{2017-2018}: Componente del Senato Accademico, University of Surrey.\\
\emph{2012-2018}: Componente del comitato per le assunzioni del personale docente del dipartimento di Economia (Job market committee member).\\
\emph{2015}: Presidente del comitato per le assunzioni del personale docente (Job market committee chair); Vice capo di dipartimento.\\
\emph{2014-2018}: Presidente della commissione per le ammissioni al Dottorato di ricerca in economia.\\
\emph{2009-2013}: Presidente della commissione per le ammissioni ai corso di laurea di primo livello del dipartimento in economia; Organizzatore dei seminari di dipartimento.\\
\\
\\
{\large \bf Competenze linguistiche e Informatiche} \\*[-.8pc]
\underline{\hspace{6in}} \\
Italiano (madrelingua), Inglese (ottima conoscenza), Spagnolo (ottima conoscenza), Portoghese (buona conoscenza);\\
Buona conoscenza di Windows, MS-Office, Mac Os e Latex. Conoscenza elementare di Linux.\\
Programmi di statistica: SPSS, Eviews, WSTATA.\\
Altri programmi: MATLAB, FORTRAN95, OX, Python, Maple e Mathematica.
\\
\\
\\
{\large \bf Attivit\`{a} di referaggio ed iscrizione a societ\`{a} scientifiche} \\*[-.8pc]
\underline{\hspace{6in}} \\
\textbf{Referee per}: \emph{Journal of Political Economy; Quarterly Journal of Economics; The Review of Economic Studies; American Economic Journal: Macroeconomics; American Economic Review: Insights; Journal of Monetary Economics; Journal of Economic Theory; Journal of the European Economic Association; The Economic Journal; Review of Economics and Statistics; Review of Economic Dynamics; Journal of Applied Econometrics; International Economic Review; Journal of Economic Behavior and Organization; Journal of Economic Dynamics and Control;  Journal of Money Credit and Banking; European Economic Review; Oxford Bulletin of Economics and Statistics; Macroeconomic Dynamics; The Scandinavian Journal of Economics; Economic Inquiry; Economic Letters; Oxford Economic papers; Journal of Macroeconomics; BE Journal of Macroeconomics; Journal of Macroeconomics; Economic Modelling; Southern Economic Journal; Scottish Journal of Political Economy; Working papers della Banca D'Inghilterra; Working papers della Deutsche Bundesbank; Working papers della Banca Nazionale del Belgio; Working papers della Banca nazionale della Repubblica Ceca; Hong Kong Research Grants Council; ESRC-UKRI; ANVUR; PRIN.}\\
\textbf{Iscritto a}: Royal economic society; American economic association; European economic association; Spanish economic association, Latin American economic association.
\\
\\
%{\large \bf Referenze} \\*[-.8pc]
%\underline{\hspace{6in}} \\
%{\bf Professor  Miguel Leon Ledesma, University of Kent.}\\
%\href{mailto:M.A.Leon-Ledesma@kent.ac.uk}{M.A.Leon-Ledesma@kent.ac.uk}\\
%{\bf Professor Peter McAdam, European Central Bank and University of Surrey.\\
%\href{mailto:peter.mcadam@ecb.int}{peter.mcadam@ecb.int}\\
%{\bf Professor Paul Levine, University of Surrey.}\\
%\href{mailto:p.levine@surrey.ac.uk}{p.levine@surrey.ac.uk}\\

%\noindent
%Informativa ai sensi dell’art.13  del D.Lgs 196/2003:\\
%i dati sopra riportati sono prescritti dalle disposizioni vigenti ai fini del procedimento per il quale sono richiesti e verranno utilizzati esclusivamente per tale scopo.
%

\bigskip
%Luogo e Data: 
Roma, 8 Giugno 2023
\bigskip


Le dichiarazioni rese nel presente curriculum sono da ritenersi rilasciate ai sensi degli artt. 46 e 47 del D.P.R. 445/2000


\begin{flushright}
Firma
\end{flushright}


%\newpage
%DICHIARAZIONI SOSTITUTIVE DI CERTIFICAZIONI (art. 46 D.P.R. n. 445/00) \\
%
%DICHIARAZIONI SOSTITUTIVE DELL’ATTO DI NOTORIETA’ (art. 47 D.P.R. n. 445/00) \\
%\\
%\\
%Il sottoscritto CANTORE CRISTIANO Codice fiscale CNTCST81E05G942N \\
%\\
%nato a POTENZA prov. PZ il 05/05/1981 sesso MASCHILE \\
%\\
%\\
%\\
%A tal fine e consapevole delle sanzioni penali, nel caso di dichiarazioni non veritiere, di formazione o uso di atti falsi, richiamate dall’art. 76 del D.P.R. 445 del 28 dicembre 2000 \\
%\\
%\\
%\begin{center}
%DICHIARA: \\
%\end{center}
%di possedere tutti i titoli riportati nel curriculum allegato alla domanda di partecipazione alla procedura di selezione relativa alla copertura di minimo n. 10 posti e massimo n. 34 posti di professore universitario di seconda fascia, da ricoprire mediante chiamata, ai sensi dell'art. 18, comma 1, della Legge 240/2010, richiesto dall'Universit`{a} degli Studi di Palermo per il CONCORSO 4 - PRIORITA IV - S.C. 13/A4 - S.S.D. SECS-P/06 - ECONOMIA APPLICATA.
%\\
%\\
%\\
%\\
%Londra, 9 Febbraio 2017\\
%\\
%\\
%\begin{center}
%il dichiarante* 
%\end{center}
%\bigskip
%\text{*}La presente dichiarazione non necessita dell’autenticazione della firma se, ai sensi dell’art. 38, D.P.R. 445/00, è sottoscritta ed inviata insieme alla fotocopia, non autenticata di un documento di identità del dichiarante, all’ufficio competente. 









\end{document} 